\section{This Thesis}

\subsection{Research Themes}
\label{sec:intro:themes}

This thesis addresses the gap between dynamics-aware and generalizable manipulation through three research themes, each targeting a different aspect of the problem.

\paragraph{RT1: Sampling-Based Kinodynamic Planning.}
The first theme develops planners that incorporate dynamics directly into the search process through physics simulation and kinodynamic exploration.
This includes introducing impulsive non-prehensile manipulation primitives, extending kinodynamic planning to settings with coupled robot-object dynamics, and developing multi-query planners that amortize dynamics computation across repeated queries.

\paragraph{RT2: Generative Models for Dynamics-Aware Motion.}
The second theme reframes motion generation as conditional sampling: a generative model learns to produce trajectories that implicitly satisfy dynamic constraints by conditioning on task-level properties such as payload mass or embodiment configuration.
This approach achieves constant-time generation without runtime constraint checking and generalizes across failure modes and payload levels.% without task-specific redesign.

\paragraph{RT3: Seeding for Constrained Trajectory Optimization.}
The third theme closes the loop between generative and optimization-based approaches by developing a task-agnostic trajectory seeding strategy.
By introducing a novel conditioning mechanism for incorporating arbitrary kinematic and dynamic constraints, this approach provides near-feasible initialization seeds that improve solver convergence and feasibility across diverse tasks without task-specific redesign.

Together, these themes trace a progression from dynamics embedded in search, to dynamics embedded in learned representations, to dynamics embedded in initialization seeds for optimization.

\subsection{Thesis Statement}
\label{sec:intro:statement}

\textit{By explicitly incorporating dynamic constraints into the motion generation process---through sampling-based kinodynamic planning, learned generative models, and dynamics-informed trajectory seeding---robots can perform a substantially broader range of manipulation tasks while operating closer to their true physical limits.}

This thesis advances robot manipulation by developing a progression of methods that expand the scope of manipulation beyond geometric pick-and-place to include dynamic non-prehensile skills and motions governed by object inertial properties, impulsive contact, and motion-induced forces.
Across all three research themes, dynamic constraints are treated not as a complexity to be avoided but as structured information to be modeled, learned, and exploited.

\subsection{Contributions}
\label{sec:intro:contributions}

The contributions of this thesis span the three research themes above.

\paragraph{Sampling-Based Kinodynamic Planning.}
\begin{enumerate}
  \item \textbf{A unifying framework for dynamics-constrained motion generation} (\cref{ch:framework}).
        A conceptual framework that organizes the landscape of planning approaches along four design axes---modeling burden, dimensionality, integration, and usability---and motivates the design choices made throughout this work.

  \item \textbf{PokeRRT: poking as a dynamic manipulation primitive} (\cref{ch:poking}).
        A sampling-based kinodynamic planner that uses physics simulation as a forward model to enable non-prehensile object rearrangement via impulsive contact, without requiring explicit contact models.
        Demonstrated across six scenarios, poking outperforms pushing and grasping baselines in success rate and task time.

  \item \textbf{Kinodynamic planning for constrained coupled-dynamics settings} (\cref{ch:constrained}).
        Two extensions of kinodynamic planning to settings with coupled robot-object dynamics: whole-body non-prehensile manipulation under locked multi-joint failures, maintaining up to 92\% of the nominal reachable area; and spill-free liquid container transport in cluttered environments, achieving a 3$\times$ expansion of the feasible solution space over prior methods.

  \item \textbf{LazyBoE: multi-query kinodynamic planning via lazy edge bundles} (\cref{ch:multiquery}).
        A multi-query kinodynamic planner that pre-computes control-parameterized edge bundles over discretized state and control spaces and applies lazy propagation and collision checking to substantially reduce repeated planning cost.
\end{enumerate}

\paragraph{Generative Models for Dynamics-Aware Motion.}
\begin{enumerate}
  \setcounter{enumi}{4}
  \item \textbf{Planning as conditional sampling} (\cref{ch:generation}).
        A reframing of dynamics-constrained motion planning as a conditional sampling problem, establishing the conceptual bridge from kinodynamic search to learned generative models.

  \item \textbf{Conditional trajectory diffusion for dynamic constraint satisfaction} (\cref{ch:payload}).
        A diffusion-based trajectory generator conditioned on payload mass that achieves super-nominal payload manipulation---67.6\% workspace accessibility at $3\times$ rated capacity---in approximately 10\,ms without runtime constraint checking.
        A case study extends this conditioning principle to fail-active manipulation under actuation failures.
\end{enumerate}

\paragraph{Seeding for Constrained Trajectory Optimization.}
\begin{enumerate}
  \setcounter{enumi}{6}
  \item \textbf{Task-agnostic seeding for constrained trajectory optimization} (\cref{ch:seeding}).
        A conditioning mechanism that generates near-feasible trajectory seeds by representing dynamic constraints as robot state samples, improving solver convergence and feasibility across diverse tasks without task-specific engineering.
\end{enumerate}
